\subsection{The cusp-swap form of the level-three gate}

Write \(\mathcal M=\operatorname{PSL}_2(\mathbf Z)\), with generators
\[
 u=\begin{pmatrix}1&1\\0&1\end{pmatrix},\qquad
 v=\begin{pmatrix}1&0\\-1&1\end{pmatrix},\qquad s=uvu.
\]
All matrix identities in this subsection are in the centreless
group. Put \(K_0=\Gamma_0(3)=\langle u,v^3\rangle\) and
let \(d_3:K_0\to\mathcal M\) be conjugation by
\(\operatorname{diag}(3,1)\). These maps also denote their
extensions to full profinite completions.

\begin{lemma}\label{lem:level3-cusp-swap-gate}
There is an involution \(\vartheta_3\) of \(K_0\) with
\[
 \vartheta_3(u)=v^3,\qquad \vartheta_3(v^3)=u,\qquad
 \vartheta_3=\operatorname{Ad}(s)\circ d_3.
\]
For an element \(g=(1,f_g)\) of full profinite GT with
\(\rho_2(g)=0\),
put \(q=f_g(u^2,v^2)\). Then
\[
 d_3(q)=q\quad\Longleftrightarrow\quad
 \vartheta_3(q)=q^{-1}.
\]
\end{lemma}

\begin{proof}
The rational matrix
\[
 W_3=\begin{pmatrix}0&1\\-3&0\end{pmatrix}
       =s\operatorname{diag}(3,1)
\]
conjugates \(u\) to \(v^3\), and \(v^3\) to \(u\).
Its square is scalar. It therefore normalizes \(K_0\) and
induces the displayed involution. The same identities hold
after full profinite completion, since they are identities
between homomorphisms of finite-index discrete groups.

Normalized congruence triviality puts \(q\) in
\(\overline{K_0}\). Conjugation by \(s\) swaps \(u\) and \(v\);
the two-cycle identity consequently gives
\[
 s q s^{-1}=f_g(v^2,u^2)=q^{-1}.
\]
Combining this with \(\vartheta_3=\operatorname{Ad}(s)d_3\)
proves the equivalence.
\end{proof}

\begin{lemma}\label{lem:level3-principal-image}
Let \(P=\Gamma(3)\), the kernel of
\(\mathcal M\to\operatorname{PSL}_2(\mathbf F_3)\). Then
\[
 d_3(P)=\Gamma^0(9),\qquad
 \vartheta_3(P)=\Gamma_0(9).
\]
In particular \(\vartheta_3\) does not preserve \(P\).
The intersection \(P\cap\Gamma_0(9)\) has index \(36\) in
\(\mathcal M\); its modular curve has genus one and six cusps, of
widths \(3,3,3,9,9,9\).
\end{lemma}

\begin{proof}
A matrix in \(P\) has, up to its common scalar sign, diagonal
entries \(1\) modulo \(3\), and off-diagonal entries divisible
by \(3\). Conjugating by \(\operatorname{diag}(3,1)\)
multiplies its upper-right entry by \(3\) and divides its
lower-left entry by \(3\). Its image therefore has upper-right
entry divisible by \(9\).
Conversely a determinant-one integral matrix with upper-right
entry divisible by \(9\) has both diagonal entries congruent
to the same sign modulo \(3\). Its inverse conjugate belongs
to \(P\). This proves the first equality; conjugation by \(s\)
gives the second. For example \(v^3\in P\) does not belong
to \(\Gamma_0(9)\).

The group \(P\) is torsion-free of index \(12\) in \(\mathcal M\).
It has four cusps, each of width \(3\): the image of \(u\)
has order \(3\) in \(\operatorname{PSL}_2(\mathbf F_3)\).
Its Euler characteristic is \(-12/6=-2\), so its compactified
modular curve has genus zero.

On \(P\), choose the representative congruent to \(I\)
modulo \(3\). Its lower-left entry divided by \(3\), reduced
modulo \(3\), is a surjective homomorphism to
\(\mathbf F_3\). Its kernel is \(P\cap\Gamma_0(9)\);
thus the latter has index three in \(P\).
For a cusp represented by \(p/q\), a positive peripheral
matrix is
\[
 I+3\begin{pmatrix}-pq&p^2\\-q^2&pq\end{pmatrix}.
\]
Its value under this homomorphism is \(-q^2\) modulo \(3\).
The cusp at infinity therefore has three unramified lifts,
of width \(3\). Each of the other three cusps has one lift
of ramification index three, of width \(9\).
There are six cusps in all. Riemann--Hurwitz for this
three-sheeted cover, branched at three points, gives
\[
 2g-2=3(-2)+3(3-1)=0.
\]
Thus \(g=1\), as claimed.
\end{proof}

\begin{lemma}\label{lem:level3-cubic-frame-principal}
For the same normalized element \(g\), let \(A_g\) be the
normalized pure cubic conjugator, so that
\[
 \Phi_g(uv)=A_g(uv)A_g^{-1},\qquad
 A_g\in\widehat F(u^2,v^2)'.
\]
Then \(A_g\in\overline{\Gamma(6)}\subseteq\overline{\Gamma(3)}\).
If also \(\rho_3^I(g)=0\), the normalized cubic identity gives
\[
 f_g(u^2,uv)\in\overline{\Gamma(6)}.
\]
\end{lemma}

\begin{proof}
Normalized congruence triviality makes the image of \(A_g\)
centralize the image of \(uv\) in
\(\operatorname{PSL}_2(\mathbf F_3)\cong A_4\).
The images of \(u^2,v^2\) generate \(A_4\), so the image of
their commutator subgroup is contained in its derived subgroup
\(V_4\). The centralizer of a nontrivial \(3\)-cycle in
\(A_4\) is its cyclic group of order three, whose intersection
with \(V_4\) is trivial. Hence \(A_g\) has trivial image. Since it already lies in the closure
of \(\Gamma(2)\), it lies in the closure of
\(\Gamma(2)\cap\Gamma(3)=\Gamma(6)\).
Intersections of finite-index subgroup closures equal the closures of
their intersections.
The last assertion follows from
\(A_g^{-1}=f_g(u^2,uv)\).
\end{proof}

These reductions leave an exact full-group question:
whether the restricted conjugator satisfies
\(\vartheta_3(q)=q^{-1}\). The principal level-three group
and its image under the cusp swap have genus-zero quotients,
but their common cover has genus one. The genus-zero
rigidity theorem therefore does not apply to that common
cover. No vanishing of the full degeneracy defect is
asserted here.

\begin{lemma}[The level-two marked cover and its regular closure]\label{lem:level3-marked-regular-closure}
Let \(F=\Gamma(2)\) and \(J=F\cap\Gamma_0(3)\).
The inclusion \(J\subset F\) has index four and its
compactified modular curve has genus zero, with cusp
widths \(2,2,2,6,6,6\).
Its regular closure over \(F\) is \(\Gamma(6)\), with
deck group \(A_4\). The latter group has index \(72\) in
\(\operatorname{PSL}_2(\mathbf Z)\), twelve cusps of width
six, and compactified genus one.
\end{lemma}
\begin{proof}
Reduction modulo three maps \(F\) onto
\(\operatorname{PSL}_2(\mathbf F_3)=A_4\): the images of
\(u^2,v^2\) generate this group. The subgroup \(J\) is the
inverse image of a cyclic subgroup of order three, namely
the stabilizer of one point of \(\mathbf P^1(\mathbf F_3)\).
Thus the coset action has degree four. Each of the three
positive peripheral elements of \(F\) maps to a nontrivial
three-cycle. Its cycle structure on the four cosets is
\(1+3\). Each of the three base cusps, of width two,
therefore has one lift of width two and one of width six.
The index of \(J\) in the modular group is \(24\), and its
six cusps give
\[
 2-2g-6=-24/6=-4,
\]
hence \(g=0\).

The core of a point stabilizer in the faithful \(A_4\)
action on four points is trivial. The core of \(J\) in
\(F\) is consequently the kernel of this reduction, which
is \(\Gamma(2)\cap\Gamma(3)=\Gamma(6)\).
Its modular index is \(6\cdot12=72\). The image of \(u\)
in the corresponding modular quotient has order six,
so all its cusp widths are six and there are twelve cusps.
Finally
\[
 2-2g-12=-72/6=-12
\]
gives \(g=1\).
\end{proof}

This calculation concerns the cover with the level-two
marking preserved. Its genus-one regular closure is different
from the genus-zero curve \(X(3)\). It does not preclude a
different quotient or filling argument, but genus-zero
rigidity cannot be applied to this regular closure.

\begin{proposition}[An explicit free-group form of the level-three gate]\label{prop:level3-free-subgroup}
Put \(F=\langle a,b\rangle=\Gamma(2)\), where \(a=u^2,b=v^2\),
and \(J=F\cap\Gamma_0(3)\). Then \(J\) has the free basis
\[
 a,\quad b^3,\quad b^2ab^{-1},\quad ba^2b^{-2},
       \quad baba^{-1}b^{-1}.
\]
The restriction of the level-three degeneracy map is the
homomorphism \(d:J\to F\) determined by
\[
\begin{array}{c|c}
 \text{generator of }J&\text{image under }d\\ \hline
 a&a^3\\
 b^3&b\\
 b^2ab^{-1}&(ba)^3\\
 ba^2b^{-2}&a^{-1}b^{-2}a^{-1}b^{-1}\\
 baba^{-1}b^{-1}&a^{-1}b^3a
\end{array}
\]
These formulas extend to full profinite completions. For
\(g=(1,f_g)\) in full GT with \(\rho_2(g)=0\), the exact
level-three gate is therefore the explicit equality
\[
             \widehat d(f_g(a,b))=f_g(a,b)
                       \quad\text{in }\widehat F.
\]
No equality in this last display is asserted by this proposition.
\end{proposition}
\begin{proof}
The left cosets of \(J\) in \(F\) are identified with the
projective lower rows of matrices modulo three. A Schreier
transversal is \(1,b,b^2,ba\), corresponding to
\([0:1],[1:1],[2:1],[1:0]\). Right multiplication by the
positive free generators gives the table
\[
\begin{array}{c|cc}
 t&\overline{ta}&\overline{tb}\\\hline
 1&1&b\\
 b&ba&b^2\\
 b^2&b&1\\
 ba&b^2&ba
\end{array}
\]
The five nontrivial Schreier words
\(t x\,\overline{tx}^{-1}\) are exactly the displayed basis.

For the signed matrices
\[
 a=\begin{pmatrix}1&2\\0&1\end{pmatrix},\qquad
 b=\begin{pmatrix}1&0\\-2&1\end{pmatrix},
\]
conjugation by \(\operatorname{diag}(3,1)\) sends
\((p,q;r,s)\) to \((p,3q;r/3,s)\).
The five basis words above have lower-left entry divisible
by three. Substitution verifies all five displayed images.
The identification with the modular degeneracy map follows
from its definition as this conjugation. The faithful
embedding of the discrete free group in the modular group
then proves the word identities.

A finite-index subgroup has its full profinite topology.
Both inclusions \(J\to F\) and the map \(d\), whose image
is the finite-index subgroup \(F\cap\Gamma^0(3)\), therefore
extend injectively to completions. Normalized congruence
triviality puts \(f_g(a,b)\) in \(\widehat J\).
The two elements in the last equality belong to
\(\widehat F\), whose inclusion in the completed modular
group is injective. Their equality is thus exactly the
previous full modular gate.
\end{proof}

The permutation action on the four cosets is the primitive
action of \(A_4\) on four points: its point stabilizer has
order three and is maximal. Consequently this particular
degree-four cover has no nontrivial intermediate cover.
It cannot be decomposed as two successive degree-two
cyclic covers. This does not rule out an identity deduced
from a larger system of covers, but the cyclic index-three
\((K_3)\) relation is a different restriction and filling map.
